\setcounter{figure}{0}
\setcounter{table}{0}
\subsection{Moral Hazard Effects using Full Unbalanced Panel}\label{subsec:app-unbal}

Our main results in \Cref{tab:mh_results} (\Cref{sec:rf}) are derived with a balanced panel to mitigate the concern of selective attrition by monitoring status. We then extend the time horizon but use the full unbalanced panel to produce \Cref{fig:rf-illus}. This section tests the robustness of both results to panel balancing.  \Cref{tab:mh_results_unbal} replicates Table \ref{tab:mh_results} in the full unbalanced panel, while \Cref{fig:rf-illus-bal} replicates \Cref{fig:rf-illus} in the balanced panel. The results presented here are almost identical to their counterparts in \Cref{sec:rf}, except for changes in power due to the differences in sample size. This suggests that our main moral-hazard results are unaffected by consumer attrition at renewals.

\begin{figure}[htbp!]
\centering
\includeifexists[scale=0.5]{\exhibitpath/appendix/fig_c1.png}
\caption{Claim Progression across Monitoring Groups (Balanced Panel) \label{fig:rf-illus-bal}}
\vspace{0.1cm}
\begin{minipage}{0.95\textwidth} % choose width suitably
    {\footnotesize \emph{Notes: }This graph reports the robustness fixed effect estimates of \cref{eq:mh_reg_long}, where we use a balanced panel for the first six periods (three years). The grey line plots $\omega_t$ while the orange line plots $\omega_t + \theta_t$, both against insurance periods $t$. The red box is superimposed ex-post to represent the period when opt-in consumers are monitored. Error-bars report 95\% confidence interval. \par}
\end{minipage}
\par
\end{figure}

\inputifexists{\exhibitpath/appendix/tab_c1}

\newpage
\subsection{Heterogeneity in Moral Hazard Effects}\label{subsec:app-mhhet}

\paragraph{Reduced-form and Model Estimates}

We investigate heterogeneity in the moral hazard effect across consumers with different observable characteristics and different insurance coverage choices. We adapt \Cref{eq:mh_reg_long} to the following specification:
\begin{align}
C_{it}= &
\alpha+\left[ \mathbf{\tau} m_{i}+ \mathbf{\omega}\mathbf{1}_{post,t} +
\mathbf{\theta_{mh}} m_{i}\cdot\mathbf{1}_{post,t} \right] \cdot (1,\mathbf{x}_{i0},y_{i0})^\prime +
(\mathbf{x}_{it},\mathbf{y}_{it})^{\prime}\mathbf{\beta} + \epsilon_{it}\label{eq:mh_reg_het}
\end{align}
Here, $(1,\mathbf{x}_{i0},y_{i0})^\prime$ indicates the initial characteristics and coverage choice of consumer $i$ and is thus time-invariant. This vector is interacted with the first and second differencing terms in  \Cref{eq:mh_reg} to uncover heterogeneity in the moral hazard effect. The interaction coefficients are presented in Figures \ref{fig:mh_reg_het_X} and \ref{fig:mh_reg_het_y}.

As the figures show, heterogeneity in the moral hazard effect is limited. The most notable variation is by drivers' initial risk class.
A positive coefficient here is intuitive: a higher risk class implies higher baseline prices and greater potential saving from the monitoring program. Drivers with higher risk class may thus be more incentivized to exert effort and reduce their accident risk.

However, such responses to incentives do not hold across other price shifters. For instance, the coefficient on zipcode income is not significant, and we do not see smaller moral hazard effects as monitoring rewards became less financially generous over the three monitoring regimes: the coefficients on state-level regime indicators (ubi\_vers\_st2 and ubi\_vers\_st3) are both positive and insignificant. This is echoed in our estimation sample, where the discount schedule shrank and the surcharge schedule grew over time, as shown in \Cref{fig:score-discount-by-regime}. However, as Table \ref{tab:app-hyperp-nonx} shows, the moral hazard effect increased across the three monitoring regimes. This suggests that the moral hazard effect might have been more heavily influenced by non-financial features of the program design (such as improved feedback mechanisms during the monitoring period) than by the scope of potential savings.

\paragraph{Model and Counterfactual Robustness}

We conduct a robustness exercise allowing the moral hazard effect to vary across risk-class in our estimation and counterfactual simulations. Tables \ref{tab:fit_c_mh_het} and \ref{tab:fit_d_mh_het} show the fit of the model, which compared to Tables \ref{tab:fit_c} and \ref{tab:fit_d}, shows no clear sign of improvement. This is unsurprising given that, when limited to the focal-state estimation data, the coefficient on risk class is no longer economically or statistically significant (\Cref{tab:app-mh-het-est}). The counterfactual results are shown in \Cref{tab:ctf_robust_mh_het}. Our main counterfactual results remain qualitatively identical and quantitatively similar.

\begin{figure}[H]
\centering
\includeifexists[scale=0.6]{\exhibitpath/appendix/fig_c2.png}
\caption{Heterogeneity in Incentive Effect across Coverage Choice \label{fig:mh_reg_het_X}}
\vspace{0.5cm}
\includeifexists[scale=0.6]{\exhibitpath/appendix/fig_c3.png}
\caption{Heterogeneity in Incentive Effect across Observables \label{fig:mh_reg_het_y}}
\vspace{0.25cm}
\begin{minipage}{0.9\textwidth}
{\footnotesize \emph{Notes:} These figures plot the estimated coefficients $\hat{\mathbf{\theta_{mh,x}}}$ in Equation (\ref{eq:mh_reg_het}) as well as the corresponding 95\% confidence intervals. A positive coefficient means that drivers with higher values (or a 1 in the case of binary variables) in the variables listed in the horizontal axis saw higher claim increase after monitoring, hence have larger moral hazard effect. The statistically significant variables are: `\text{x\_drvr\_hm\_own\_ind}' (home ownership indicator), `rc' (risk class), `cov\_ded\_coll\_ind' (driver has collision coverage ind.), `cov\_ind\_fullcov' (driver has both collision and comprehensive coverage ind.). \par}
\end{minipage}
\end{figure}

\inputifexists{\exhibitpath/appendix/tab_c2}
\inputifexists{\exhibitpath/appendix/tab_c3}
\inputifexists{\exhibitpath/appendix/tab_c4}

\pagebreak
\subsection{Learning Effects from Monitoring}\label{subsec:app-learning}

Our baseline specification assumes that monitoring does not alter drivers' persistent risk types. In other words, we do not model learning effects whereby drivers permanently reduce their risk after being monitored. Allowing for such learning would not change our moral hazard estimates, but it could reduce the extent of advantageous selection, since part of the post-monitoring risk gap between monitored and unmonitored groups could then reflect learning rather than selection.

\paragraph{Evidence from Public Records of At-Fault Accidents}

Because monitoring begins in the first policy period, the firm's internal claims data cannot capture any pre-existing risk differences between monitored and unmonitored drivers. To assess this, we turn to drivers' public-record at-fault accidents (AFA records). For all new customers, the firm purchases a LexisNexis report containing a historical aggregate measure of at-fault accidents, covering one to five years of prior records depending on driver experience and state reporting rules. This measure enters the firm's risk-class assignment and thus pricing. For renewals, the firm does not purchase new reports; instead, it reports actuarially adjusted claims to LexisNexis and internally updates AFA records accordingly.

Compared to claims data, AFA records are delayed when claims involve complex adjustments that spill over into future periods. Moreover, some claims, especially minor ones, may be missing from AFA records due to administrative leniency or incomplete adjustments. Overall, the mean trailing-twelve-month AFA records for a renewal customer is 0.018, compared to the average claim records of 0.05 per six-month period.

However, these limitations are not systematically affected by monitoring. We can therefore replicate the difference-in-differences analysis in Section \ref{eq:mh_reg_long}, replacing claims with the occurrence of at-fault accidents from the public record. We find that the monitored group has 8.7\% lower AFA in the pre-period (the level gap shown in \Cref{fig:rf-mh-viol} divided by the mean AFA in the unmonitored group of the corresponding period), which changed to 38.8\% during monitoring, 21.9\% and 7.4\% post-monitoring.

Consistent with our main results, monitored drivers appear safer even in the pre-period, indicating the presence of advantageous selection. The moral hazard effect is evident during monitoring, when AFA records further diverge between monitored and unmonitored groups, but this effect completely disappears when we re-examine the AFA records between one to two years after monitoring (at the beginning of year 3). We thus conclude that monitoring's learning effect, as measured by its persistent impact on accident risk, is likely negligible.

\begin{figure}[H]
\centering
\includeifexists[scale=0.75]{\exhibitpath/appendix/fig_c4.png}
\caption{At-Fault Accident Violation Progression by Monitoring Groups \label{fig:rf-mh-viol}}
\begin{minipage}{0.95\textwidth} % choose width suitably
    {\footnotesize \emph{Notes: }
    This graph reports the year fixed-effect estimates across monitoring groups of a revised version of \cref{eq:mh_reg_long}, where we use public-record at-fault accidents as the dependent variable. For renewal customers, we focus on the beginning of policy periods 1 (right after monitoring), 3 (year 1 post-monitoring), and 5 (year 2 post-monitoring). For new customers, historical aggregates cover one to five years depending on experience and reporting rules. Risk class is excluded from controls; results remain similar when it is included. The grey line plots $\omega_t$ and the orange line $\omega_t + \theta_t$. Error bars denote 95\% confidence intervals.\par}
\end{minipage}
\end{figure}

\paragraph{Counterfactual Robustness}

As a robustness check, we re-estimate the model and counterfactuals under the assumption of a 10\% learning effect, i.e., that monitored drivers' accident risk falls permanently by 10\% after monitoring.\footnote{We choose 10\% because, in our analysis above, relative to the pre-period, the change in the monitored-unmonitored AFA gap is -13.2\% and +1.2\% one and two years after monitoring. Even if we completely ignore the reporting delay in AFA records and the last data point, 10\% would represent the larger end of learning effects consistent with the data.} Table \ref{tab:ctf_robust_learning} shows that learning increases both consumer and producer surplus, leading to more aggressive monitoring pricing in the initial period and greater rent-sharing in renewals. However, our main results---the welfare benefits from monitoring and the harm of data portability regulation---remain unchanged.

\inputifexists{\exhibitpath/appendix/tab_c5}

\vspace{-1em}
